Rozdział drugi poświęcony jest szczegółowej analizie algorytmów zaproponowanych w pracy Fornasiera, Schnassa i Vybírala \cite{fornasier2012learning}. Głównym problemem rozważanym przez autorów jest efektywne uczenie (aproksymacja) funkcji wielowymiarowej zdefiniowanej na jednostkowej kuli w przestrzeni $\mathbb{R}^d$, przyjmującej postać $f(x)=g(Ax)$. W modelu tym $A$ jest macierzą o wymiarach $k \times d$, a $g$ jest funkcją $k$ zmiennych. Kluczowym założeniem jest to, że wymiar $k$ jest znacznie mniejszy od $d$ ($k \ll d$). Klasyczne metody numeryczne w obliczu tak postawionego zadania natrafiają na problem praktycznej nierozwiązywalności (tzw. klątwy wymiarowości), ponieważ liczba próbek wymagana do aproksymacji rośnie wykładniczo wraz z wymiarem $d$.Zamiarem autorów było stworzenie algorytmów, których złożoność rośnie co najwyżej wielomianowo względem wymiaru $d$ oraz liczby punktów pomiarowych $m$. Osiągnięto to dzięki umiejętnemu wykorzystaniu faktu, że jeśli funkcja zależy jedynie od kilku kombinacji liniowych zmiennych wejściowych, to jej gradienty leżą w niskowymiarowej podprzestrzeni i wykazują właściwości wektorów rzadkich lub ściśliwych. Umożliwia to zastosowanie aparatu teorii próbkowania skompresowanego (ang. compressed sensing) do stabilnego odzyskania macierzy parametrów z relatywnie niewielkiej liczby losowych zapytań o wartości funkcji. Wyprowadzenie algorytmów zaproponowanych przez Vybírala i współautorów opiera się na przybliżeniu pochodnych kierunkowych za pomocą rozwinięcia Taylora. Dla punktu próbkowania $\xi \in B_{\mathbb{R}^d}$ oraz kierunku $\varphi \in B_{\mathbb{R}^d}(r)$, i dla odpowiednio małego parametru $\epsilon \in \mathbb{R}_+$, zachodzi następująca tożsamość:


$$[\nabla g(A\xi)^T A]\varphi = \frac{f(\xi+\epsilon\varphi)-f(\xi)}{\epsilon} - \frac{\epsilon}{2}[\varphi^T \nabla^2 f(\zeta) \varphi],$$

gdzie $\zeta$ jest odpowiednio dobranym punktem pośrednim z otoczenia.Powtarzając takie przybliżenie dla $m_\Phi$ losowych kierunków (tworzących macierz $\Phi$) oraz $m_\mathcal{X}$ punktów próbkowania (tworzących zbiór $\mathcal{X}$), algorytmy sprowadzają problem identyfikacji parametrów do układu równań:$$\Phi X = Y - \mathcal{E}$$gdzie $Y$ to macierz różnic skończonych (estymat pochodnych), $\mathcal{E}$ to błąd przybliżenia wzoru Taylora, a nieznana macierz $X$ ma specyficzną strukturę $X = A^T \mathcal{G}^T$ (dla macierzy gradientów $\mathcal{G}$). Kolumny macierzy $X$ stanowią kombinacje liniowe wierszy macierzy $A$, przez co są ściśliwe.Z uwagi na to, że układ ten reprezentuje problem z dziedziny compressed sensing, w pierwszym kroku obu algorytmów kolumny macierzy $X$ są estymowane niezależnie za pomocą minimalizacji normy $l_1$, dając przybliżenie $\hat{X}.$

Dalsze procedury zależą od wymiaru parametru:
\begin{enumerate}
    \item W przypadku jednowymiarowym ($k=1$, Algorytm 1): Macierz $A$ redukuje się do pojedynczego wektora $a$. Wektory odzyskane za pomocą minimalizacji $l_1$ są przeskalowanymi kopiami wektora $a^T$. Algorytm wyznacza ostateczne przybliżenie $\hat{a}$ po prostu jako znormalizowany wektor $\hat{x}_j$ posiadający największą normę $l_2$ (co odpowiada punktowi o największym gradiencie).
    \item W przypadku wielowymiarowym ($k \ge 1$, Algorytm 2): Struktura jest bardziej złożona. Aby wyekstrahować poprawną bazę z macierzy przybliżeń $\hat{X}$, algorytm stosuje rozkład według wartości osobliwych (SVD) dla macierzy $\hat{X}^T$. Ponieważ prawdziwa macierz $X$ ma rząd nie większy niż $k$, przybliżenie macierzy $\hat{A}$ jest konstruowana z $k$ prawych wektorów osobliwych odpowiadających największym wartościom osobliwym macierzy $\hat{X}^T$.
\end{enumerate}
